\documentclass[11pt]{article}
\usepackage[margin=1in]{geometry}
\usepackage[T1]{fontenc}
\usepackage{lmodern}
\usepackage{amsmath,amssymb,amsthm}
\usepackage[hidelinks]{hyperref}
\usepackage{microtype}

\newtheorem{theorem}{Theorem}
\newtheorem{lemma}[theorem]{Lemma}
\newtheorem{corollary}[theorem]{Corollary}
\newcommand{\U}{\mathcal U}
\newcommand{\N}{\mathbb N}

\begin{document}

\begin{center}
{\Large Full-result packet for arXiv:1605.06909}\\[4pt]
{\large A contractible, unitarily representable SIN Polish group\\
that is not of finite type}
\end{center}

\paragraph{Status.}
\textbf{full\_solution\_likely\_valid}.  This packet answers Question~4.5
of Hiroshi Ando, Yasumichi Matsuzawa, Andreas Thom, and Asger T\"ornquist,
\emph{Unitarizability, Maurey--Nikishin factorization, and Polish groups of
finite type}, arXiv:1605.06909v3.

\paragraph{Source question.}
A Polish group is \emph{unitarily representable} if it is topologically
isomorphic to a subgroup of the unitary group of a separable Hilbert space,
with the strong operator topology.  It is \emph{of finite type} if it is
topologically isomorphic to a closed subgroup of the unitary group of a
finite von Neumann algebra with separable predual.  The source constructs
unitarily representable SIN Polish groups which are not of finite type, but
its construction has a discrete quotient.  It then asks:

\begin{quote}
Is there a connected, unitarily representable SIN Polish group which is not
of finite type?
\end{quote}

We give an affirmative answer with ``connected'' strengthened to
``contractible.''  The construction is a general permanence device.

\paragraph{The measurable-function group.}
Let \(G\) be a Polish group and let

\[
 L^0(G)=L^0([0,1],G)
\]

be the group of equivalence classes, modulo equality almost everywhere, of
Lebesgue-measurable maps \(f:[0,1]\to G\), with pointwise multiplication and
the topology of convergence in measure.  If \(d\leq1\) is a compatible
metric on \(G\), that topology is induced by

\[
 D_d(f,h)=\int_0^1 d(f(x),h(x))\,dx.
\]

It is standard, and also follows by approximating with simple functions,
that \(L^0(G)\) is Polish when \(G\) is Polish.

\begin{theorem}[The \(L^0\)-permanence device]\label{thm:device}
Let \(G\) be a unitarily representable SIN Polish group.  Then \(L^0(G)\)
is a contractible, unitarily representable SIN Polish group.  Moreover, if
\(G\) is not of finite type, then \(L^0(G)\) is not of finite type.
\end{theorem}

\begin{proof}
We prove the four assertions separately.

\smallskip
\noindent\emph{SIN and Polishness.}
Choose a bounded compatible two-sided invariant metric \(d\) on \(G\).  The
integral metric \(D_d\) above is two-sided invariant:

\[
 D_d(af,ah)=D_d(f,h)=D_d(fa,ha)
 \qquad(a,f,h\in L^0(G)).
\]

It induces convergence in measure.  Thus \(L^0(G)\) is SIN, and the standard
simple-function completeness and separability argument makes it Polish.

\smallskip
\noindent\emph{Contractibility.}
For \(0\leq s\leq1\), define

\[
 [C_s f](x)=
 \begin{cases}
 f(x),&0\leq x\leq1-s,\\
 e,&1-s<x\leq1,
 \end{cases}
\]

where \(e\) is the identity of \(G\).  Then \(C_0f=f\) and \(C_1f=e\).
For all \(s,t\in[0,1]\) and \(f,h\in L^0(G)\), boundedness of \(d\) gives

\[
 D_d(C_sf,C_th)\leq D_d(f,h)+|s-t|.
\]

Hence \((s,f)\mapsto C_sf\) is a continuous contraction of \(L^0(G)\) to
its identity.  In particular, \(L^0(G)\) is path connected and connected.

\smallskip
\noindent\emph{Unitary representability.}
Let

\[
 \rho:G\longrightarrow\U(H)
\]

be a topological group embedding, where \(H\) is separable and
\(\U(H)\) has the strong operator topology.  On
\(K=L^2([0,1],H)\), define the multiplication representation

\[
 [\Pi(f)\xi](x)=\rho(f(x))\xi(x).
\]

This is a well-defined injective homomorphism \(L^0(G)\to\U(K)\).  We check
that it is a topological embedding, which is the only point requiring care.

Choose a dense sequence \((v_j)_{j\geq1}\) in the unit sphere of \(H\).  The
bounded metric

\[
 r(g,h)=\sum_{j=1}^{\infty}2^{-j}
 \min\bigl\{1,\|[\rho(g)-\rho(h)]v_j\|\bigr\}
\]

induces the topology of \(G\).  Consequently \(D_r\) induces convergence in
measure on \(L^0(G)\).

Suppose \(D_r(f_n,f)\to0\).  For every \(j\), the bounded functions

\[
 x\longmapsto\|[\rho(f_n(x))-\rho(f(x))]v_j\|
\]

converge to zero in measure, hence in \(L^2\).  Thus
\(\Pi(f_n)\to\Pi(f)\) strongly on every constant vector field \(v_j\).
Finite-valued fields with values in the span of the \(v_j\)'s are dense in
\(K\), and all the operators are unitary, so the convergence is strong on
all of \(K\).

Conversely, if \(\Pi(f_n)\to\Pi(f)\) strongly, testing on each constant
field \(v_j\) gives

\[
 \int_0^1\|[\rho(f_n(x))-\rho(f(x))]v_j\|^2\,dx\longrightarrow0.
\]

The defining series for \(r\), first truncated and then estimated by its
uniformly small tail, yields \(D_r(f_n,f)\to0\).  Therefore \(\Pi\) is a
topological embedding into the unitary group of the separable Hilbert space
\(K\).

\smallskip
\noindent\emph{Failure of finite type.}
The constant-map embedding

\[
 \iota:G\longrightarrow L^0(G),\qquad
 \iota(g)(x)=g,
\]

is a topological group embedding with closed image.  To see closedness,
if a sequence of constant maps converges in measure, pass to a subsequence
which converges almost everywhere.  At any two points in the common
conull set, the same sequence in \(G\) has both pointwise limits, so those
limits agree; the limiting function is constant almost everywhere.

Finite type passes to closed subgroups directly from its definition.  Thus,
if \(L^0(G)\) were of finite type, its closed constant subgroup
\(\iota(G)\) would be of finite type.  This proves the contrapositive of the
last assertion.
\end{proof}

\begin{corollary}[Answer to the source question]
There exists a contractible, unitarily representable SIN Polish group which
is not of finite type.
\end{corollary}

\begin{proof}
Choose a uniformly bounded non-unitarizable representation
\(\pi:\mathbb F_2\to\mathrm{GL}(H)\) on a separable Hilbert space and form

\[
 G_0=H\rtimes_{\pi}\mathbb F_2.
\]

The Main Theorem of arXiv:1605.06909 says that \(G_0\) is a unitarily
representable SIN Polish group and that it is of finite type if and only if
\(\pi\) is unitarizable.  Hence \(G_0\) is not of finite type.  Apply
Theorem~\ref{thm:device} to obtain \(L^0(G_0)\).
\end{proof}

\paragraph{Checks and scope.}
The argument uses the source's definitions of unitary representability and
finite type.  The multiplication representation is not merely faithful: the
constant-vector-field calculation proves that it induces exactly the
convergence-in-measure topology.  The non-finite-type conclusion does not
require a new von Neumann algebra obstruction; it is inherited from the
closed subgroup of constant functions.  Contractibility is asserted only as
a property of the underlying topological space, which is more than the
connectedness requested by the source.

\paragraph{Novelty check.}
Before promotion, the run indexes were searched for arXiv:1605.06909,
``connected unitarily representable SIN,'' ``Polish groups of finite type,''
and measurable-function groups; no earlier packet was found.  The current
arXiv record remains v3 and contains the question.  Bounded searches for the
exact sentence and for \(L^0([0,1],G)\) together with finite type found the
source and general work on measurable-function groups, but no stated answer
by this permanence device.

\paragraph{Human-review recommendation.}
Review as a short full answer.  The two points worth checking are the reverse
topology implication for the multiplication representation and the use of
closed-subgroup heredity for finite type; both are written out above.

\begin{thebibliography}{9}
\bibitem{source}
H. Ando, Y. Matsuzawa, A. Thom, and A. T\"ornquist,
\emph{Unitarizability, Maurey--Nikishin factorization, and Polish groups of
finite type}, arXiv:1605.06909v3, 2017.

\bibitem{l0}
A. Kwiatkowska and M. Malicki,
\emph{Automorphism groups of countable structures and groups of measurable
functions}, arXiv:1612.03106, 2016.
\end{thebibliography}

\end{document}
